\chapter{Literature Review}
\label{chap:lit}

\section{Review scope and evaluation framework}

The proposed system combines four fields that are often treated separately: pavement and ride-quality sensing, embedded signal processing, vehicular telemetry, and automotive Head-Up Display (HUD) design. The literature is therefore evaluated against the needs of the TUKZIE Rev 0 rather than only by reported classification accuracy. Relevant evidence must also address computational cost, sampling requirements, sensitivity to vehicle speed and mounting position, real-time latency, communication loss, and driver workload. This framework is important because a technique that performs well on a passenger car, smartphone, or laboratory simulator may not transfer directly to a lightweight three-wheeled cargo platform.

\section{Light electric freight vehicles and last-mile logistics}

Light electric freight vehicles can reduce the energy, space, and emissions associated with urban delivery, but their advantage depends on route length, payload, consolidation strategy, electricity supply, and the vehicle used for comparison. In Dar es Salaam, Sebwa et al. compared an electric cargo tricycle with motorcycles, petrol tuk-tuks, and a light-duty car on selected medium-size delivery routes \cite{sebwa2025}. The electric tricycle had the lowest reported operating cost, with savings ranging from 45.5\% against motorcycles to 86\% against the light-duty car. The study also reported energy use of up to 0.54 kWh for a one-way trip and a 97.59\% energy saving relative to the light-duty-car comparison \cite{sebwa2025}. It avoided the direct tailpipe emissions produced by the petrol vehicles. These results provide valuable regional evidence, although the underlying fuel prices, route choices, assumed payload, trip length, comparison basis, and omission of electricity-generation emissions limit direct transfer to Cape Town.

Evidence from Portland reaches a similar conclusion using a different system boundary. Saenz et al. combined operational GPS and warehouse data with vehicle, battery, fuel, and electricity life-cycle contributions, finding greenhouse-gas reductions of 51--72\% for an electric-tricycle logistics service relative to conventional urban logistics \cite{saenz2016}. The range is more informative than a single headline value: it shows that benefits change when partner operations and supporting logistics are included. For TUKZIE, the literature therefore justifies the electric cargo-trike context but does not by itself demonstrate that one prototype will reduce costs or emissions. The present project contributes enabling telemetry that can later measure route energy, utilisation, ride exposure, and operational reliability on the actual platform.

Three-wheel geometry and a compact chassis also make vehicle response an engineering concern rather than merely a logistics detail. Changes in payload position, tyre pressure, speed, and road surface can affect both stability and measured vibration. This strengthens the case for recording operating state alongside IMU measurements and for avoiding the direct transfer of passenger-car roughness thresholds to the TUKZIE platform.

\section{Pavement roughness and vehicle ride smoothness}

Road roughness and ride smoothness are related but not interchangeable concepts. The International Roughness Index (IRI) is calculated from a measured longitudinal road profile using a standard quarter-car simulation \cite{fhwa2016}. It provides a transferable engineering measure when the profile measurement and calibration procedures are controlled. By contrast, an accelerometer mounted on a vehicle measures the combined response of the road input, tyres, suspension, chassis, payload, speed, sensor orientation, and mounting location. Consequently, raw vertical acceleration cannot be interpreted as IRI without a validated mapping and suitable reference measurements.

Gonz\'alez et al. demonstrated that road profiles can be classified from simulated axle and body accelerations, establishing the value of vehicle-response measurements for roughness estimation \cite{gonzalez2008}. The approach is attractive because it reuses vehicle dynamics rather than requiring a dedicated laser profiler. Its limitation is equally important: the inferred roughness depends on the assumed vehicle model and operating conditions. A model calibrated for a conventional four-wheeled vehicle cannot be transferred uncritically to the TUKZIE trike, whose wheel layout, suspension response, payload distribution, and low operating speed differ.

Smartphone and low-cost inertial sensing studies extend this response-based approach. Sattar et al. found that accelerometers and positioning sensors can support scalable road surface monitoring, but identified device placement, orientation, speed, sampling rate, and inconsistent ground truth as persistent sources of variation \cite{sattar2018}. Alatoom and Smadi addressed cost and portability through an embedded framework using microcontroller and single-board computer platforms \cite{alatoom2025}. Their architecture is closer to the present project than a smartphone-only system because it permits controlled sensor mounting, deterministic sampling, and local processing. Nevertheless, published low-cost systems generally target passenger vehicles or generic road-anomaly detection; they do not establish a calibrated ride-smoothness measure for a three-wheeled electric cargo vehicle.

These findings support a deliberately bounded claim for this project. The TUKZIE unit should first produce a repeatable, vehicle-specific ride-smoothness index and event log. Correlation with IRI may be investigated when reference profile data are available, but the system should not label its output as IRI solely from uncalibrated chassis acceleration. Repeatability over repeated traversals, sensitivity to speed and payload, and agreement with a reference sensor or labelled route provide more defensible validation criteria.

\section{Inertial signal processing and vibration separation}
\label{sec:lit-signal-processing}

An IMU signal contains gravitational acceleration, orientation changes, road excitation, vehicle-body motion, motor and drivetrain vibration, sensor noise, and transient events. The processing pipeline must therefore preserve road-related information while removing components that would make comparisons between runs misleading. Orientation compensation or transformation into a vehicle-fixed frame is required before features from the nominally vertical axis are compared. Resampling and timestamp checks are also necessary because feature values become unreliable when samples are irregular or missing.

The literature uses both time-domain and frequency-domain features. Root-mean-square acceleration, standard deviation, peak-to-peak amplitude, crest factor, jerk, and threshold counts are computationally inexpensive and interpretable. Frequency-domain energy and power spectral density can distinguish sustained narrow-band vibration from broadband road impacts. Bui et al. show that signal-processing-driven rough-road detection can be executed on an edge microcontroller, supporting the feasibility of local classification without continuous cloud processing \cite{bui2026}. However, threshold-based methods can overfit a particular vehicle and speed, while more complex classifiers can hide the physical reason for a detection and demand a larger labelled dataset.

Recent smartphone research illustrates both the opportunity and this limitation. Yiliguoqi et al. used vertical acceleration and roll-derived features, including standard deviation and interquartile range, with an SVM at a 10 Hz sampling rate to distinguish roughness classes \cite{yiliguoqi2026}. The low sampling frequency and compact feature set are attractive for constrained hardware. The reported experiments classified the defined road-quality levels correctly in 80--100\% of trials across four vehicles and repeated short road segments \cite{yiliguoqi2026}. However, that range is tied to the authors' phones, routes, labels, IRI thresholds, and limited test segments; it cannot be treated as an expected accuracy for the TUKZIE without retraining and independent validation. SVM or XGBoost models should therefore be considered as comparison methods after a transparent statistical baseline has been established, not as substitutes for calibration.

For the TUKZIE platform, separation of road excitation from motor and powertrain vibration requires an empirical baseline. Stationary tests at several motor speeds can identify repeatable spectral components associated with the powertrain. Coast-down and repeated route tests can then show which components persist without powered traction and which vary with the road surface. Windowed statistics should be normalised or stratified by vehicle speed, since the same spatial road feature produces a different temporal excitation as speed changes. Rather than relying on one feature, a small interpretable feature vector combining vertical RMS acceleration, jerk, robust peak measures, and energy in selected frequency bands offers a defensible balance between sensitivity and computational cost.

Whole-body vibration standards provide useful terminology but do not replace road-profile measurement. ISO 2631-1 defines methods for evaluating human exposure to whole-body vibration \cite{iso2631}. Its frequency-weighted measures can inform a comfort-oriented index, whereas IRI remains a road-profile statistic. Reporting these as separate outputs would prevent the common error of treating passenger comfort, vehicle response, and pavement geometry as the same quantity.

\section{Edge computing and real-time telemetry architecture}

Edge computing moves computation closer to the sensors to reduce response time, bandwidth use, and dependence on a remote service \cite{shi2016}. These benefits directly match the TUKZIE use case: the HUD and safety alerts must continue to operate when cellular coverage is poor, and transmitting every raw high-frequency sample would consume more bandwidth and energy than transmitting compact features and events. The embedded HPEC concept also enables several sensor streams to be timestamped and fused before data leave the vehicle \cite{barnell2018}.

The term ``high performance'' must nevertheless be interpreted relative to the platform. The relevant performance is not maximum throughput, but predictable completion of time-sensitive work within a constrained power and thermal budget. A multithreaded design can isolate high-rate sensor acquisition from feature extraction, HUD rendering, local storage, and networking, as required by the project brief \cite{winberg2026}. Threads alone do not guarantee real-time behaviour: shared locks, blocking input/output, memory allocation, and unbounded queues can still cause missed samples or stale displays.

A producer--consumer architecture is therefore appropriate. The acquisition task should timestamp samples and place them into bounded buffers; processing tasks should consume complete windows; the HUD should read the most recent validated state without waiting for network activity; and the communication task should transmit queued summaries asynchronously. Overload behaviour must be explicit. Dropping an old display frame may be acceptable, whereas silently losing an IMU window or corrupting timestamps is not. Performance should be reported as distributions of acquisition jitter, processing time, display age, queue occupancy, and end-to-end latency rather than as an average alone.

P\'erez-Gonz\'alez et al. provide a useful contemporary comparison through a field-tested electric-racing telemetry system that combines ESP32 acquisition, Raspberry Pi 4B edge analytics, Node-RED visualisation, and LoRaWAN communication \cite{perezgonzalez2025}. Its modular division between deterministic acquisition and higher-level processing supports the architecture proposed here. The racing study reports a 99\% packet-delivery ratio and sub-second telemetry latency during its field trials \cite{perezgonzalez2025}. These results demonstrate feasibility for that track, gateway placement, message rate, and radio configuration; they do not establish LoRaWAN as universally superior in moving vehicles. LoRaWAN offers low-power long-range packets when gateways are available, whereas 4G/LTE provides wider public-network reach and direct Internet connectivity at the cost of subscription, coverage dependence, and less predictable delay. The comparison supports retaining local logging and asynchronous communication regardless of which radio is selected.

\section{HUD design, attention, and visual workload}

HUDs reduce the need for large downward eye movements, but their position in the forward field of view does not make them automatically safe. Liu reported that HUD use changes attention demand and driving performance, showing that display location and cognitive processing must be evaluated together \cite{liu2003}. Smith et al. similarly found that conventional head-down-display assessment methods do not fully capture HUD behaviour and that both driving and secondary-task performance should be considered \cite{smith2016}. The design problem is therefore not simply to project more telemetry onto the windshield.

Windshield-display research identifies a design space involving location, visual depth, contrast, information density, and the relationship between graphics and the external scene \cite{hauslschmid2016}. Peripheral presentation can support awareness, but misaligned or excessive graphics may compete with hazards for attention \cite{hauslschmid2015}. Optical quality also depends on the windshield and projection system; distortion, double images, ambient illumination, and viewing position affect legibility \cite{mareska2022}. These constraints are especially relevant to a removable, low-cost projector mounted on a trike, where vibration and changing daylight are likely to be more severe than in an integrated production vehicle.

Ghosting arises because the inner and outer windshield surfaces can produce spatially separated reflections of the same projected content. Wedge-shaped glass or PVB interlayers can reduce this separation for a defined optical geometry, but manufacturing tolerances and movement away from the design viewing position prevent a universal correction \cite{mareska2022, qin2017}. The appropriate acceptance criterion is therefore empirical legibility across the intended viewing region rather than a wedge angle copied from another vehicle. For example, Qin et al. calculated approximately 0.62 mrad for overlap of the primary and ghost images in their particular model \cite{qin2017}; this is a case-specific result rather than a TUKZIE design value. The eye box, the region in which the complete virtual image remains visible, must likewise be measured for the TUKZIE installation. Published automotive systems use different eye-box dimensions according to their field of view and optical architecture, so a nominal dimension is a design benchmark rather than a guaranteed requirement.

Complexity is a central risk. Currano et al. showed that HUD complexity influences driver perceptions of automated vehicles \cite{currano2021}, while cognitive-load experiments on connected HUDs demonstrate that adding screens or information channels can increase mental demand \cite{zhu2021}. The TUKZIE HUD should therefore use a stable information hierarchy: persistent speed and battery state, brief prioritised safety alerts, and navigation cues only when required. Decorative graphics, rapidly changing values, and simultaneous low-priority notifications should be excluded. The system should also adopt fail-safe behaviour when data are stale, for example by removing or clearly marking invalid values rather than continuing to display them as current.

Park and Im found symbol count to be the most influential HUD-interface factor in their semi-automated-driving study and recommended displaying fewer than six symbols simultaneously \cite{park2020}. This is a useful initial ceiling for the prototype, not a universal safety threshold: the meaning, salience, location, duration, and driving context of the symbols remain important. Combined with the NHTSA two-second glance guidance \cite{nhtsa2013}, it supports a testable design requirement for a sparse, prioritised interface.

Evaluation should combine technical and human-centred measures. Technical measures include rendering time, update rate, luminance or contrast under representative lighting, and the age of displayed telemetry. Human-centred evaluation can use task completion, missed-event counts, glance behaviour, perceived workload, and structured feedback during controlled low-risk trials. This is more defensible than asserting that a HUD is safer solely because it is head-up.

\section{Cellular synchronisation and data integrity}

Cellular communication extends local telemetry to fleet and maintenance applications, but coverage and delay are outside the embedded controller's direct control. The architecture should consequently separate the real-time control path from remote synchronisation. HUD operation and local event detection must not wait for a 4G/LTE response. Each record should carry a monotonic sample timestamp, a route or session identifier, sensor-quality flags, and a sequence number so that the server can detect gaps and duplicates.

MQTT is suitable for compact publish--subscribe telemetry, whereas HTTP is useful for configuration, bulk upload, and conventional database interfaces. Protocol choice alone does not ensure reliability. A store-and-forward queue, bounded retry policy, acknowledgement strategy, and reconnection test are required. The appropriate performance measures are local data completeness, queued backlog, synchronisation delay, duplicate rate, and recovery after an induced outage. These measures directly address the project's remote data-synchronisation deliverable \cite{winberg2026}.

\section{Proximity sensing on low-speed vehicles}

Live nearby-object alerting is a recognised requirement even for low-speed vehicles, not only high-speed autonomous ones. De Simone et al. used ultrasonic range data, classified by a supervised neural network, to add obstacle avoidance to retrofitted autonomous agricultural machines \cite{desimone2018}, showing that low-cost range sensing can support obstacle awareness on slow-moving vehicles. Closer to this project's own vehicle class, Noomwongs et al. developed a low-speed autonomous micro electric vehicle that combines RTK-GNSS localisation with a 2D LiDAR for real-time obstacle avoidance, explicitly because vision-based approaches demand expensive sensors and computing power that such vehicles lack, a constraint this project shares \cite{noomwongs2025}, confirming that proximity/obstacle awareness is treated as a genuine requirement for this vehicle category in the literature, not only for highway-speed autonomous cars.

This project's own live-alerting requirement is addressed with a time-of-flight (ToF) rather than an ultrasonic or vision-based sensor, and a solid-state multi-zone ToF sensor (the VL53L5CX) was selected over a single-point alternative specifically for its ability to resolve direction (left/centre/right) from a single device, discussed further in Chapter~\ref{chap:method}. An independent laboratory characterisation of this exact sensor \cite{caroleo2026} is used directly in that chapter to set realistic, evidence-based expectations for its accuracy and reliability, rather than relying on the manufacturer's datasheet figures alone.

\section{Synthesis and research gap}

Existing studies establish that low-cost inertial sensors can identify road anomalies, edge devices can process sensor windows locally, and HUDs can reduce downward glances. They also expose three limitations. First, acceleration-based roughness estimates are strongly dependent on the vehicle, speed, mounting, and calibration. Second, edge architectures are often evaluated by algorithm accuracy or mean latency without reporting jitter, queue behaviour, and outage recovery. Third, HUD benefits diminish when visual complexity, optical limitations, and cognitive workload are ignored.

No reviewed work integrates these concerns for a low-cost, removable telemetry and windshield-HUD module designed specifically for the vibration, payload variation, limited on-board resources, and intermittent connectivity of a lightweight electric cargo trike in an emerging-market context. High-end passenger vehicles integrate telemetry and optically engineered HUDs, while low-cost road-sensing and IoV studies usually investigate their subsystems separately. This project addresses that gap by combining controlled IMU mounting, vehicle-specific vibration baselining, interpretable edge features, non-blocking local processing, store-and-forward telemetry, and a minimal HUD information hierarchy. The contribution is therefore not merely the coexistence of a sensor and display. It is the design and validation of an integrated architecture whose ride-smoothness repeatability, latency, data integrity, and presentation safety can be measured explicitly.
